\section{Build, Deployment, and HIL Setup}\label{sec:deployment}

\subsection{Linux Daemon}

The Linux reference binary is \texttt{thermalcored}. \texttt{make build} produces it under \texttt{build/platform-linux/}; it loads a JSON configuration, binds platform backends (\texttt{hwmon} or file sensors, \texttt{sysfs} or HIL PWM, SocketCAN context), starts the configured telemetry transport, and calls \texttt{thermal\_core\_step()} once per control period:

\begin{lstlisting}[caption={Running the Linux daemon}]
make build
build/platform-linux/thermalcored --config=my-config.json
\end{lstlisting}

The same daemon serves all three Linux roles --- a standalone Raspberry~Pi~4 regulator, the scenario-clock host for the determinism harness, and the host side of the hardware-in-the-loop setup --- selected entirely by configuration.

\subsection{ESP32-C3 Firmware: Three Build Modes}

The reference firmware targets the ESP32-C3 and links the same portable \texttt{core/} and \texttt{protocol/} C99 sources as the Linux daemon. It builds in three mutually exclusive modes, selected by a CMake flag (the build system rejects any attempt to enable two at once):

\begin{table}[h]
\centering
\begin{tabular}{p{0.26\textwidth}p{0.62\textwidth}}
\toprule
\textbf{Mode} & \textbf{Behavior} \\
\midrule
\texttt{STANDALONE} (default) & The ESP32 reads its DS18B20 probe and fan tachometer, runs \texttt{thermal\_core\_step()} locally, drives the fan over LEDC PWM, and prints a human-readable status line. \\
\texttt{REPLAY\_STANDALONE} & Synthetic input snapshots baked into \texttt{.rodata} replace the live sensors; the core runs locally and emits a byte-stable CSV. Used for on-target determinism checks without any wiring. \\
\texttt{HIL\_PERIPHERAL} & The ESP32 owns only the peripherals; the Linux daemon runs the core. See below. \\
\bottomrule
\end{tabular}
\caption{ESP32-C3 firmware build modes.}
\label{tab:esp32-modes}
\end{table}

\begin{lstlisting}[caption={Building and flashing the firmware}]
cd platform/esp32_idf
idf.py build                                  # STANDALONE
idf.py -DTHERMALCORE_REPLAY_STANDALONE=ON build
idf.py -DTHERMALCORE_HIL_PERIPHERAL=ON  build
idf.py flash monitor                           # flash + serial console

# all three modes + size-budget check, from the repo root:
make build-esp32
\end{lstlisting}

\subsection{Hardware-in-the-Loop}

\texttt{HIL\_PERIPHERAL} mode splits the system across the USB cable. The ESP32-C3 owns the board-support layer --- DS18B20 over 1-Wire, the tachometer GPIO interrupt, LEDC PWM --- and nothing else. It samples the peripherals, encodes each reading as a \texttt{TELEM\_SAMPLE} \texttt{TC} frame, and streams it over USB-Serial-JTAG to the Linux daemon. The daemon runs \texttt{thermal\_core\_step()}, then sends the resulting fan duty back as a \texttt{CMD\_REQUEST} frame, which the firmware applies to the PWM output. The control law executes on the host while the real sensor and the real fan stay on the bench.

Because both ends link the identical \texttt{core/} source, HIL is also a portability check: a control decision computed by the Linux daemon over HIL is the decision the \texttt{STANDALONE} firmware would have computed on the same inputs. A bench bring-up of this loop is reported in Section~\ref{sec:evaluation}.

\subsection{OBD-II Emulator}

\texttt{car-can-emulator} is included as a submodule under \texttt{tools/car-can-emulator/} and tracks the upstream \texttt{v2-improvements} branch. It is a test dependency, not a runtime dependency of the thermal core: it synthesizes OBD-II vehicle-speed responses on a virtual CAN interface so the acoustic-mask policy and the CAN fault paths can be exercised in CI without a vehicle.

\begin{lstlisting}[caption={Emulator build}]
git submodule update --init --recursive
cmake -H tools/car-can-emulator -B build/car-can-emulator
cmake --build build/car-can-emulator
\end{lstlisting}

\subsection{Step-by-Step Guides}

The repository carries three task-oriented guides under \texttt{docs/getting-started/} that walk through each deployment with exact wiring and commands: \texttt{esp32-standalone.md} (a self-contained ESP32-C3 regulator), \texttt{linux-pi4.md} (a Raspberry~Pi~4 daemon driving a \texttt{sysfs} PWM fan), and \texttt{hil-peripheral.md} (the ESP32-plus-daemon HIL loop). This section gives the architecture; those guides give the keystrokes.

\endinput
